\clearpage
\begin{appendices}
\renewcommand{\thechapter}{\Asbuk{chapter}}
\titleformat{\chapter}[display]{\filcenter}{\MakeUppercase{Приложение} \thechapter}{0pt}{\bfseries}

\chapter{КОНТРОЛЬНЫЕ МАТЕРИАЛЫ И ФРАГМЕНТЫ КОДА}

\section{Контрольные точки и методика верификации}

Верификация вычислительного ядра выполнена на основе контрольных материалов IAPWS-IF97.  
Контрольные данные включают набор точек по различным регионам стандарта и проверяют корректность вычисления базовых свойств (например, $v$, $h$, $s$) при заданных параметрах $p$ и $T$, а также корректность обратных маршрутов в пределах области применимости \cite{iapws_if97_2012}.  

Для автоматизации проверки применяются модульные и интеграционные тесты.  
Модульные тесты проверяют отдельные формулы и вспомогательные функции (например, определение региона и работу с границами области).  
Интеграционные тесты проверяют маршруты расчёта на уровне публичного фасада \texttt{If97} и сверяют результаты с эталонными значениями в допустимой погрешности.  

Отдельное внимание уделяется особым зонам области применимости.  
К таким зонам относятся линия насыщения (Region~4) и околокритическая область (Region~3), где свойства резко меняются, а обратные зависимости требуют устойчивых численных процедур.  

\section{Пример использования API вычислительного ядра (Rust)}

Ниже приведён укрупнённый пример вызовов публичного API ядра \texttt{if97\_core}.  
В реальном проекте в интерфейсных оболочках дополнительно выполняется разбор пользовательского ввода, валидация формата данных и отображение результатов.  

\begin{verbatim}
use if97_core::{If97, units::{MegaPascal, Kelvin}};

fn main() -> Result<(), if97_core::domain::If97Error> {
    let if97 = If97::new();

    // pT: direct calculation
    let state = if97.pt(MegaPascal(10.0), Kelvin(773.15))?;
    println!("h = {} kJ/kg, s = {} kJ/(kg*K)", state.h, state.s);

    // pH: backward calculation
    let state2 = if97.ph(MegaPascal(15.0), state.h)?;
    println!("T = {} K, region = {}", state2.t, state2.region);

    Ok(())
}
\end{verbatim}

\section{Фрагмент DTO для обмена между слоями}

Для взаимодействия между UI-частью и backend-слоем используется отдельный крейт \texttt{if97\_app\_api}.  
Он содержит сериализуемые структуры, не зависящие от внутренних типов ядра, и обеспечивает стабильность контрактов обмена при развитии приложения.  
В частности, в DTO учитываются ограничения сериализации (например, необходимость корректного кодирования специальных числовых значений).  

\chapter{ДОПОЛНИТЕЛЬНЫЕ СХЕМЫ, ДИАГРАММЫ И МАНИФЕСТЫ}

\section{Схема потоков данных в программном комплексе}

В проекте реализован подход «одно ядро --- несколько потребителей», в котором \texttt{if97\_core} остаётся единственным источником вычислений.  
Укрупнённая схема потоков данных приведена ниже.  

\begin{verbatim}
                 +------------------+
                 |     if97_core     |
                 | (single compute)  |
                 +---------+---------+
                           |
         +-----------------+------------------+
         |                 |                  |
 +-------+------+   +------+--------+   +-----+------+
 |  FLTK UI     |   |  Yew/WASM UI  |   | gRPC service|
 | (desktop)    |   | + Tauri backend|   | (batch API) |
 +--------------+   +---------------+   +-------------+
\end{verbatim}

\section{Фрагмент gRPC-протокола и SoA-представление данных}

Для сервисного контура применяется gRPC \cite{grpc_docs}.  
Ниже приведён укрупнённый пример описания сообщений, использующих SoA-представление данных.  
Такой формат снижает накладные расходы при пакетной обработке и упрощает data-parallel вычисления.  

\begin{verbatim}
message PtBatchRequest {
  uint64 batch_id = 1;
  repeated double p = 2; // MPa
  repeated double t = 3; // K
}

message PtBatchResponse {
  uint64 batch_id = 1;
  repeated double h = 2; // kJ/kg
  repeated uint32 status = 3;
}
\end{verbatim}

\section{Перечень параметров конфигурации сервиса}

Конфигурация сервисного слоя выполняется через переменные окружения, что упрощает развертывание в контейнерной и оркестраторной среде.  
К основным параметрам относятся:
\begin{itemize}
\item \texttt{IF97\_GRPC\_ADDR} --- адрес и порт gRPC-сервера;  
\item \texttt{IF97\_IO\_THREADS} и \texttt{IF97\_CPU\_THREADS} --- настройка параллелизма I/O и CPU;  
\item \texttt{IF97\_MAX\_IN\_FLIGHT\_PER\_CONN} и \texttt{IF97\_MAX\_IN\_FLIGHT\_GLOBAL} --- лимиты для backpressure;  
\item \texttt{IF97\_MAX\_BATCH\_LEN} и \texttt{IF97\_GRPC\_MAX\_MSG\_BYTES} --- ограничения на размеры батчей и сообщений;  
\item \texttt{IF97\_DRAIN\_SECS} --- время корректного завершения (drain) при остановке.  
\end{itemize}

\section{Фрагменты инфраструктурных файлов (Docker, Kubernetes, CI/CD)}

Контейнеризация сервиса выполнена с использованием Docker \cite{docker_docs}.  
Эксплуатационная инфраструктура ориентирована на развёртывание в Kubernetes \cite{kubernetes_docs} и на автоматизацию сборки и релизов через GitHub Actions \cite{github_actions_docs}.  

Ниже приведены укрупнённые примеры фрагментов конфигураций, используемых в проекте.  

\textbf{Kubernetes Deployment (фрагмент):}
\begin{verbatim}
env:
  - name: IF97_GRPC_ADDR
    value: "0.0.0.0:50051"
  - name: IF97_CPU_THREADS
    value: "8"
\end{verbatim}

\textbf{GitHub Actions workflow (фрагмент):}
\begin{verbatim}
on:
  push:
    tags:
      - "v*.*.*"
\end{verbatim}

\end{appendices}
